\documentclass[11pt]{article}
\usepackage[a4paper,margin=22mm]{geometry}
\usepackage{fontspec,polyglossia}
\setmainlanguage{english}\setotherlanguages{italian,latin}
\setmainfont{Linux Libertine O}[Ligatures=TeX]
\setlength{\parindent}{0pt}
\begin{document}
\begin{center}{\Large Boncompagni's Italian in his edition of \textit{Algoritmi de numero Indorum} (1857)}\\[1mm]
{\large in English}\end{center}
{\small Baldassarre Boncompagni printed the Latin arithmetic «Dixit Algorizmi» from the Cambridge manuscript Ii.\,6.\,5
in the first volume of his \textit{Trattati d'aritmetica} (Rome, 1857). His own matter is in Italian: the title page,
the heading of the text, the folio marks of the manuscript in the margin and six footnotes. Each line is given as this
repository transcribes it (\texttt{../anchors.tsv}), with its English; the Latin he quotes stands in italics with an
English rendering in parentheses. The Latin text of the treatise is not translated here: J. N. Crossley and A. S.
Henry translated the Cambridge manuscript into English («Thus Spake al-Khwārizmī», \textit{Historia Mathematica} 17,
1990). AI-translated (Claude), 2026; no person has proofread it. Dedicated to the public domain (CC0 1.0).}
\medskip
\subsection*{Title page}
\noindent\begin{tabular}{@{}p{0.2\textwidth}@{\hspace{2mm}}p{0.38\textwidth}@{\hspace{3mm}}p{0.38\textwidth}@{}}{\footnotesize l. 1} & {\small TRATTATI} & TREATISES\end{tabular}\par\smallskip
\noindent\begin{tabular}{@{}p{0.2\textwidth}@{\hspace{2mm}}p{0.38\textwidth}@{\hspace{3mm}}p{0.38\textwidth}@{}}{\footnotesize l. 2} & {\small D’ ARITMETICA} & ON ARITHMETIC\end{tabular}\par\smallskip
\noindent\begin{tabular}{@{}p{0.2\textwidth}@{\hspace{2mm}}p{0.38\textwidth}@{\hspace{3mm}}p{0.38\textwidth}@{}}{\footnotesize l. 3} & {\small PUBBLICATI} & PUBLISHED\end{tabular}\par\smallskip
\noindent\begin{tabular}{@{}p{0.2\textwidth}@{\hspace{2mm}}p{0.38\textwidth}@{\hspace{3mm}}p{0.38\textwidth}@{}}{\footnotesize l. 4} & {\small DA BALDASSARRE BONCOMPAGNI} & BY BALDASSARRE BONCOMPAGNI\end{tabular}\par\smallskip
\noindent\begin{tabular}{@{}p{0.2\textwidth}@{\hspace{2mm}}p{0.38\textwidth}@{\hspace{3mm}}p{0.38\textwidth}@{}}{\footnotesize l. 5} & {\small \textsc{socio ordinario dell’accademia pontificia de’nuovi lincei,}} & ordinary member of the Pontifical Academy of the Nuovi Lincei,\end{tabular}\par\smallskip
\noindent\begin{tabular}{@{}p{0.2\textwidth}@{\hspace{2mm}}p{0.38\textwidth}@{\hspace{3mm}}p{0.38\textwidth}@{}}{\footnotesize l. 6} & {\small \textsc{e socio corrispondente dell’accademia reale}} & and corresponding member of the Royal Academy\end{tabular}\par\smallskip
\noindent\begin{tabular}{@{}p{0.2\textwidth}@{\hspace{2mm}}p{0.38\textwidth}@{\hspace{3mm}}p{0.38\textwidth}@{}}{\footnotesize l. 7} & {\small \textsc{delle scienze di torino,}} & of Sciences of Turin,\end{tabular}\par\smallskip
\noindent\begin{tabular}{@{}p{0.2\textwidth}@{\hspace{2mm}}p{0.38\textwidth}@{\hspace{3mm}}p{0.38\textwidth}@{}}{\footnotesize l. 8} & {\small \textsc{della reale accademia delle scienze di napoli,}} & of the Royal Academy of Sciences of Naples,\end{tabular}\par\smallskip
\noindent\begin{tabular}{@{}p{0.2\textwidth}@{\hspace{2mm}}p{0.38\textwidth}@{\hspace{3mm}}p{0.38\textwidth}@{}}{\footnotesize l. 9} & {\small \textsc{e della pontificia accademia delle scienze}} & and of the Pontifical Academy of Sciences\end{tabular}\par\smallskip
\noindent\begin{tabular}{@{}p{0.2\textwidth}@{\hspace{2mm}}p{0.38\textwidth}@{\hspace{3mm}}p{0.38\textwidth}@{}}{\footnotesize l. 10} & {\small \textsc{dell’ istituto di bologna}} & of the Institute of Bologna\end{tabular}\par\smallskip
\noindent\begin{tabular}{@{}p{0.2\textwidth}@{\hspace{2mm}}p{0.38\textwidth}@{\hspace{3mm}}p{0.38\textwidth}@{}}{\footnotesize l. 11} & {\small I.} & I.\end{tabular}\par\smallskip
\noindent\begin{tabular}{@{}p{0.2\textwidth}@{\hspace{2mm}}p{0.38\textwidth}@{\hspace{3mm}}p{0.38\textwidth}@{}}{\footnotesize l. 12} & {\small ALGORITMI DE NUMERO INDORUM.} & ALGORITMI DE NUMERO INDORUM. (Algoritmi on the numbers of the Indians.)\end{tabular}\par\smallskip
\noindent\begin{tabular}{@{}p{0.2\textwidth}@{\hspace{2mm}}p{0.38\textwidth}@{\hspace{3mm}}p{0.38\textwidth}@{}}{\footnotesize l. 13} & {\small ROMA} & ROME\end{tabular}\par\smallskip
\noindent\begin{tabular}{@{}p{0.2\textwidth}@{\hspace{2mm}}p{0.38\textwidth}@{\hspace{3mm}}p{0.38\textwidth}@{}}{\footnotesize l. 14} & {\small TIPOGRAFIA DELLE SCIENZE FISICHE E MATEMATICHE} & PRESS OF THE PHYSICAL AND MATHEMATICAL SCIENCES\end{tabular}\par\smallskip
\noindent\begin{tabular}{@{}p{0.2\textwidth}@{\hspace{2mm}}p{0.38\textwidth}@{\hspace{3mm}}p{0.38\textwidth}@{}}{\footnotesize l. 15} & {\small in Via Lata N.° 211.} & at Via Lata no. 211.\end{tabular}\par\smallskip
\noindent\begin{tabular}{@{}p{0.2\textwidth}@{\hspace{2mm}}p{0.38\textwidth}@{\hspace{3mm}}p{0.38\textwidth}@{}}{\footnotesize l. 16} & {\small 1857.} & 1857.\end{tabular}\par\smallskip
\subsection*{Heading of the text}
\noindent\begin{tabular}{@{}p{0.2\textwidth}@{\hspace{2mm}}p{0.38\textwidth}@{\hspace{3mm}}p{0.38\textwidth}@{}}{\footnotesize p. 1, l. 1} & {\small I.} & I.\end{tabular}\par\smallskip
\noindent\begin{tabular}{@{}p{0.2\textwidth}@{\hspace{2mm}}p{0.38\textwidth}@{\hspace{3mm}}p{0.38\textwidth}@{}}{\footnotesize p. 1, l. 2} & {\small \textbf{ALGORITMI DE NUMERO INDORUM.}} & ALGORITMI DE NUMERO INDORUM. (Algoritmi on the numbers of the Indians.)\end{tabular}\par\smallskip
\noindent\begin{tabular}{@{}p{0.2\textwidth}@{\hspace{2mm}}p{0.38\textwidth}@{\hspace{3mm}}p{0.38\textwidth}@{}}{\footnotesize p. 1, l. 3} & {\small (Biblioteca dell’Università di Cambridge, Codice manoscritto contrassegnato « \textit{Ii}. 6. 5 » carte 102, \textit{recto} — 109, \textit{verso}).} & (University Library, Cambridge, manuscript marked «Ii. 6. 5», ff. 102 recto – 109 verso.)\end{tabular}\par\smallskip
\subsection*{The folio marks in the margin and the footnotes}
\noindent\begin{tabular}{@{}p{0.2\textwidth}@{\hspace{2mm}}p{0.38\textwidth}@{\hspace{3mm}}p{0.38\textwidth}@{}}{\footnotesize p. 1, margin, at l. 5} & {\small car. 102 \textit{recto}.} & f. 102 recto.\end{tabular}\par\smallskip
\noindent\begin{tabular}{@{}p{0.2\textwidth}@{\hspace{2mm}}p{0.38\textwidth}@{\hspace{3mm}}p{0.38\textwidth}@{}}{\footnotesize p. 1, footnote} & {\small (1) La lacuna indicata con quattro punti nelle linea 25 della presente pagina trovasi anche nel sopraccitato Codice \textit{Ii}. 6. 5 della Biblioteca dell’ Università di Cambridge (carta 102 \textit{recto}, lin. 18).} & (1) The lacuna marked by four dots in line 25 of this page is found also in the above-mentioned manuscript Ii. 6. 5 of the University Library, Cambridge (f. 102 recto, l. 18).\end{tabular}\par\smallskip
\noindent\begin{tabular}{@{}p{0.2\textwidth}@{\hspace{2mm}}p{0.38\textwidth}@{\hspace{3mm}}p{0.38\textwidth}@{}}{\footnotesize p. 2, footnote} & {\small (1) La lacuna indicata con quattro punti nella linea terza della presente pagina trovasi anche nel sopraccitato Codice \textit{Ii}. 6. 5. (carta 102 \textit{recto}, lin. 22).} & (1) The lacuna marked by four dots in the third line of this page is found also in the above-mentioned manuscript Ii. 6. 5 (f. 102 recto, l. 22).\end{tabular}\par\smallskip
\noindent\begin{tabular}{@{}p{0.2\textwidth}@{\hspace{2mm}}p{0.38\textwidth}@{\hspace{3mm}}p{0.38\textwidth}@{}}{\footnotesize p. 2, margin, at l. 17} & {\small car. 102 \textit{verso}.} & f. 102 verso.\end{tabular}\par\smallskip
\noindent\begin{tabular}{@{}p{0.2\textwidth}@{\hspace{2mm}}p{0.38\textwidth}@{\hspace{3mm}}p{0.38\textwidth}@{}}{\footnotesize p. 4, margin, at l. 1} & {\small car. 103 \textit{recto}.} & f. 103 recto.\end{tabular}\par\smallskip
\noindent\begin{tabular}{@{}p{0.2\textwidth}@{\hspace{2mm}}p{0.38\textwidth}@{\hspace{3mm}}p{0.38\textwidth}@{}}{\footnotesize p. 4, footnote} & {\small (1) Le due lacune, ciascuna delle quali è indicata con quattro punti nella prima linea della presente pagina, trovansi anche nel sopraccitato Codice \textit{Ii}. 6. 5. (carta 102 \textit{verso}, lin. 33).} & (1) The two lacunae, each of which is marked by four dots in the first line of this page, are found also in the above-mentioned manuscript Ii. 6. 5 (f. 102 verso, l. 33).\end{tabular}\par\smallskip
\noindent\begin{tabular}{@{}p{0.2\textwidth}@{\hspace{2mm}}p{0.38\textwidth}@{\hspace{3mm}}p{0.38\textwidth}@{}}{\footnotesize p. 5, margin, at l. 23} & {\small car. 103 \textit{verso}.} & f. 103 verso.\end{tabular}\par\smallskip
\noindent\begin{tabular}{@{}p{0.2\textwidth}@{\hspace{2mm}}p{0.38\textwidth}@{\hspace{3mm}}p{0.38\textwidth}@{}}{\footnotesize p. 7, margin, at l. 6} & {\small car. 104 \textit{recto}.} & f. 104 recto.\end{tabular}\par\smallskip
\noindent\begin{tabular}{@{}p{0.2\textwidth}@{\hspace{2mm}}p{0.38\textwidth}@{\hspace{3mm}}p{0.38\textwidth}@{}}{\footnotesize p. 7, footnote} & {\small (1) La lacuna indicata con quattro punti nella linea 17 della presente pagina trovasi anche nel sopraccitato Codice \textit{Ii}. 6. 5. (carta 104 \textit{recto}, lin. 8, 9).} & (1) The lacuna marked by four dots in line 17 of this page is found also in the above-mentioned manuscript Ii. 6. 5 (f. 104 recto, ll. 8, 9).\end{tabular}\par\smallskip
\noindent\begin{tabular}{@{}p{0.2\textwidth}@{\hspace{2mm}}p{0.38\textwidth}@{\hspace{3mm}}p{0.38\textwidth}@{}}{\footnotesize p. 8, margin, at l. 17} & {\small car. 104 \textit{verso}.} & f. 104 verso.\end{tabular}\par\smallskip
\noindent\begin{tabular}{@{}p{0.2\textwidth}@{\hspace{2mm}}p{0.38\textwidth}@{\hspace{3mm}}p{0.38\textwidth}@{}}{\footnotesize p. 9, margin, at l. 28} & {\small car. 105 \textit{recto}.} & f. 105 recto.\end{tabular}\par\smallskip
\noindent\begin{tabular}{@{}p{0.2\textwidth}@{\hspace{2mm}}p{0.38\textwidth}@{\hspace{3mm}}p{0.38\textwidth}@{}}{\footnotesize p. 9, footnote} & {\small (1) La lacuna indicata con quattro punti nella linea 29 della presente pagina trovasi anche nel sopraccitato Codice \textit{Ii}. 6. 5. (carta 105 \textit{recto} , linea prima).} & (1) The lacuna marked by four dots in line 29 of this page is found also in the above-mentioned manuscript Ii. 6. 5 (f. 105 recto, first line).\end{tabular}\par\smallskip
\noindent\begin{tabular}{@{}p{0.2\textwidth}@{\hspace{2mm}}p{0.38\textwidth}@{\hspace{3mm}}p{0.38\textwidth}@{}}{\footnotesize p. 11, margin, at l. 8} & {\small car. 105 \textit{verso}.} & f. 105 verso.\end{tabular}\par\smallskip
\noindent\begin{tabular}{@{}p{0.2\textwidth}@{\hspace{2mm}}p{0.38\textwidth}@{\hspace{3mm}}p{0.38\textwidth}@{}}{\footnotesize p. 12, margin, at l. 19} & {\small car. 106 \textit{recto}.} & f. 106 recto.\end{tabular}\par\smallskip
\noindent\begin{tabular}{@{}p{0.2\textwidth}@{\hspace{2mm}}p{0.38\textwidth}@{\hspace{3mm}}p{0.38\textwidth}@{}}{\footnotesize p. 13, margin, at l. 25} & {\small car. 106 \textit{verso}.} & f. 106 verso.\end{tabular}\par\smallskip
\noindent\begin{tabular}{@{}p{0.2\textwidth}@{\hspace{2mm}}p{0.38\textwidth}@{\hspace{3mm}}p{0.38\textwidth}@{}}{\footnotesize p. 15, margin, at l. 3} & {\small car. 107 \textit{recto}.} & f. 107 recto.\end{tabular}\par\smallskip
\noindent\begin{tabular}{@{}p{0.2\textwidth}@{\hspace{2mm}}p{0.38\textwidth}@{\hspace{3mm}}p{0.38\textwidth}@{}}{\footnotesize p. 16, margin, at l. 16} & {\small car. 107 \textit{verso}.} & f. 107 verso.\end{tabular}\par\smallskip
\noindent\begin{tabular}{@{}p{0.2\textwidth}@{\hspace{2mm}}p{0.38\textwidth}@{\hspace{3mm}}p{0.38\textwidth}@{}}{\footnotesize p. 17, margin, at l. 24} & {\small car. 108 \textit{recto}.} & f. 108 recto.\end{tabular}\par\smallskip
\noindent\begin{tabular}{@{}p{0.2\textwidth}@{\hspace{2mm}}p{0.38\textwidth}@{\hspace{3mm}}p{0.38\textwidth}@{}}{\footnotesize p. 18, margin, at l. 30} & {\small car. 108 \textit{verso}.} & f. 108 verso.\end{tabular}\par\smallskip
\noindent\begin{tabular}{@{}p{0.2\textwidth}@{\hspace{2mm}}p{0.38\textwidth}@{\hspace{3mm}}p{0.38\textwidth}@{}}{\footnotesize p. 20, margin, at l. 15} & {\small car. 109 \textit{recto}.} & f. 109 recto.\end{tabular}\par\smallskip
\noindent\begin{tabular}{@{}p{0.2\textwidth}@{\hspace{2mm}}p{0.38\textwidth}@{\hspace{3mm}}p{0.38\textwidth}@{}}{\footnotesize p. 22, margin, at l. 5} & {\small car. 109 \textit{verso}.} & f. 109 verso.\end{tabular}\par\smallskip
\noindent\begin{tabular}{@{}p{0.2\textwidth}@{\hspace{2mm}}p{0.38\textwidth}@{\hspace{3mm}}p{0.38\textwidth}@{}}{\footnotesize p. 23, footnote} & {\small (1) Nel margine inferiore della carta 109 \textit{verso} del suddetto codice \textit{Ii}. 6. 5. si legge: \textit{dicque constitues} .\textsc{viii}. Più sotto nel margine stesso si legge: \textit{est quidem}. Il \textit{recto} della carta 110 del codice medesimo incomincia così (lin. 1—4): \textit{Est quidem guazal cora tabula ad celi rotunditatem conformata, cuius itemque} \textit{superficies equali spatiorum intersticione quadrifidis ab ipso centro per diametrum, fusis lineationibus, partita est.}} & (1) In the lower margin of f. 109 verso of the said manuscript Ii. 6. 5 one reads: \textit{dicque constitues .viii.} (and say: you will set down 8). Lower down in the same margin one reads: \textit{est quidem} (there is indeed). The recto of f. 110 of the same manuscript begins thus (ll. 1–4): \textit{Est quidem guazal cora tabula ad celi rotunditatem conformata, cuius itemque superficies equali spatiorum intersticione quadrifidis ab ipso centro per diametrum, fusis lineationibus, partita est.} (There is indeed a «guazal cora», a board shaped to the roundness of the sky, whose surface likewise is divided from the centre itself along the diameter into four parts, by lines drawn out at equal intervals.)\end{tabular}\par\smallskip
\end{document}
